\chapter{常用不等式}
\label{app:inequalities}

本附录整理经常用于估计、比较和误差控制的不等式，并列出其主要依据。选择不等式时，应注意当前需要控制的量以及等号对应的结构。

\section{Cauchy不等式}

对实数\(a_1,\ldots,a_n\)、\(b_1,\ldots,b_n\)，
\begin{equation*}
 \left(\sum_{i=1}^n a_ib_i\right)^2
 \le\left(\sum_{i=1}^n a_i^2\right)\left(\sum_{i=1}^n b_i^2\right).
\end{equation*}
等号当且仅当两个向量\((a_i)\)、\((b_i)\)线性相关。可由恒等式
\begin{equation*}
 \left(\sum a_i^2\right)\left(\sum b_i^2\right)-\left(\sum a_ib_i\right)^2
 =\sum_{i<j}(a_ib_j-a_jb_i)^2
\end{equation*}
直接得到。内积形式为\(|\langle u,v\rangle|\le\|u\|\|v\|\)，从而
\begin{equation*}
 \|u+v\|\le\|u\|+\|v\|,\qquad
 \bigl|\|u\|-\|v\|\bigr|\le\|u-v\|.
\end{equation*}

对\(x_i>0\)，取\(a_i=\sqrt{x_i}\)、\(b_i=1/\sqrt{x_i}\)，可得
\begin{equation*}
 \left(\sum_{i=1}^n x_i\right)\left(\sum_{i=1}^n\frac1{x_i}\right)\ge n^2.
\end{equation*}

\section{平均值不等式}

对正数\(a_1,\ldots,a_n\)，调和平均、几何平均、算术平均与平方平均满足
\begin{equation*}
 \frac{n}{\sum_{i=1}^n1/a_i}
 \le\left(\prod_{i=1}^n a_i\right)^{1/n}
 \le\frac1n\sum_{i=1}^n a_i
 \le\sqrt{\frac1n\sum_{i=1}^n a_i^2}.
\end{equation*}
各处等号均当且仅当全部\(a_i\)相等。

加权形式中，设\(w_i>0\)、\(\sum w_i=1\)，则
\begin{equation*}
 \prod_{i=1}^n a_i^{w_i}\le\sum_{i=1}^n w_i a_i.
\end{equation*}
取\(A=\sum w_i a_i\)，对每个\(a_i/A\)使用\(\ln x\le x-1\)，加权相加便得\(\sum w_i\ln(a_i/A)\le0\)，即上式。对\(1/a_i\)应用同一结论得到前面的调和平均界，平方平均界则由Cauchy不等式得到。

\section{指数、对数与幂}

常用的切线估计为
\begin{equation*}
 e^x\ge1+x\quad(x\in\mathbb R),\qquad
 \ln x\le x-1\quad(x>0).
\end{equation*}
等号分别在\(x=0\)、\(x=1\)处取得。由对数的积分表达或作差求导，\(x>0\)时还有
\begin{equation*}
 \frac{x}{1+x}<\ln(1+x)<x.
\end{equation*}

Bernoulli不等式为
\begin{equation*}
 (1+x)^n\ge1+nx\qquad(x\ge-1,\ n\in\mathbb N^*),
\end{equation*}
可以直接归纳证明。对于实数指数\(r\)与\(x>-1\)，由函数的凸性可得
\begin{equation*}
 \begin{cases}
 (1+x)^r\ge1+rx,&r\le0\text{或}r\ge1,\\
 (1+x)^r\le1+rx,&0\le r\le1.
 \end{cases}
\end{equation*}

\section{函数增量与积分估计}

若\(f\)在连接\(x,y\)的闭区间连续、内部可导，且\(|f'|\le M\)，则中值定理给出
\begin{equation*}
 |f(x)-f(y)|\le M|x-y|.
\end{equation*}
若\(f\)在\([a,b]\)可积，\(m\le f\le M\)，则
\begin{equation*}
 m(b-a)\le\int_a^bf(x)\,dx\le M(b-a),\qquad
 \left|\int_a^bf\right|\le\int_a^b|f|.
\end{equation*}
特别地，连续函数满足
\begin{equation*}
 \left|\int_a^b(f-g)\right|\le(b-a)\|f-g\|_\infty.
\end{equation*}
这些估计分别把局部变化率和整个区间上的函数误差，转化为可以计算的数值界。
